\subsection{Алгоритм ICP и регистрация облаков точек}
\label{sec:1-2-icp}

\subsubsection{Итерационная схема алгоритма ICP}

Одним из наиболее распространённых методов регистрации облаков точек является алгоритм ICP. В задаче лидарной одометрии он используется для оценки относительного движения по двум соседним кадрам, когда текущее облако точек необходимо совместить с опорным представлением сцены. Результатом работы алгоритма является евклидово преобразование, включающее поворот и перенос, которое наилучшим образом согласует геометрию двух облаков. Именно это преобразование затем интерпретируется как оценка относительного перемещения сенсора между соседними моментами времени.

Если бы соответствия между точками двух облаков были известны заранее, задача оценки преобразования сводилась бы к сравнительно простой геометрической подзадаче. Однако на практике такие соответствия неизвестны: для каждой точки текущего облака необходимо ещё определить, какая точка опорного облака описывает тот же фрагмент сцены. Тем самым задача регистрации приобретает взаимозависимый характер: чтобы корректно оценить преобразование, нужны хорошие соответствия, а чтобы найти хорошие соответствия, уже требуется достаточно точное преобразование. Именно эта зависимость и делает итерационный подход естественным решением.

Основная идея ICP состоит в попеременном выполнении двух шагов. На первом шаге при фиксированном текущем положении облака строятся соответствия между точками текущего и опорного облаков. На втором шаге при фиксированном множестве соответствий пересчитывается евклидово преобразование, которое лучше согласует найденные пары точек. После этого текущее облако переносится в новое положение, поиск соответствий выполняется заново уже для уточнённого положения, и цикл повторяется. Такой процесс продолжается до тех пор, пока изменение преобразования между соседними итерациями не станет достаточно малым, либо пока не будет достигнуто заранее заданное ограничение на число итераций. Таким образом, ICP можно рассматривать как процедуру последовательного уточнения оценки движения, в которой соответствия и преобразование улучшаются совместно от итерации к итерации; следовательно, этот алгоритм является локальным итерационным методом, и его поведение существенно зависит от начального приближения.

\subsubsection{Построение соответствий и оценка преобразования}

Наиболее простой вариант построения соответствий основан на критерии ближайшего соседа. Для каждой точки текущего облака выбирается ближайшая к ней точка опорного облака, после чего полученная пара рассматривается как кандидат на соответствие. Такой подход удобен вычислительно и хорошо работает в ситуации, когда межкадровое смещение невелико. Для ускорения поиска обычно используются специальные пространственные структуры данных, например KD-деревья~\cite{kd-tree-habr}, позволяющие быстро находить ближайшие точки в трёхмерном пространстве. Вместе с тем не каждая найденная пара должна автоматически приниматься в дальнейшую обработку: если расстояние между точками слишком велико, такое соответствие с высокой вероятностью является ошибочным и должно быть отброшено по порогу допустимого расстояния.

После построения множества соответствий требуется оценить евклидово преобразование, которое наилучшим образом совмещает найденные пары точек. В классическом point-to-point ICP эта задача обычно формулируется как минимизация суммы квадратов расстояний между соответствующими точками двух облаков. Иначе говоря, требуется подобрать такой поворот и такой перенос, чтобы после преобразования текущего облака суммарное рассогласование между всеми принятыми парами оказалось минимальным. Для этой подзадачи существуют эффективные вычислительные процедуры; в частности, широко используется сингулярное разложение матрицы (SVD)~\cite{svd-yandex-handbook}, позволяющее по уже найденным парам точек вычислить параметры преобразования.

Помимо point-to-point постановки, в литературе часто рассматривается и вариант point-to-plane ICP~\cite{point-cloud-registration-survey-huang}, в котором минимизируется не полное евклидово расстояние между точками, а расстояние от точки одного облака до локальной касательной плоскости в другом облаке. Такой вариант может обеспечивать более быструю сходимость на гладких поверхностях, однако требует дополнительной оценки нормалей и тем самым усложняет вычислительную схему. В данной работе в качестве базового метода рассматривается именно point-to-point ICP, поскольку он проще в реализации, не требует отдельного этапа вычисления нормалей и позволяет сосредоточиться на влиянии динамического перекрытия поля зрения на саму процедуру регистрации.

\subsubsection{Применимость и ограничения ICP}

Корректная работа классического ICP фактически опирается на несколько предпосылок. Во-первых, межкадровое смещение должно оставаться достаточно малым, чтобы поиск ближайших соседей находил геометрически осмысленные пары точек. Во-вторых, между соседними кадрами должно сохраняться заметное перекрытие статической части сцены. В-третьих, доля ложных соответствий не должна быть слишком велика, иначе оценка преобразования начинает смещаться. Наконец, оставшаяся наблюдаемая геометрия должна быть достаточно информативной, чтобы по ней можно было надёжно восстановить поворот и перенос.

Практическая привлекательность ICP объясняется несколькими причинами. Этот алгоритм имеет ясную геометрическую интерпретацию, непосредственно работает с облаками точек без необходимости вводить сложные промежуточные представления сцены и хорошо сочетается с последовательной обработкой лидарных кадров. При небольшом относительном смещении между соседними кадрами ICP часто обеспечивает достаточно точную оценку движения даже в базовой конфигурации. Именно поэтому он остаётся естественным базовым методом для задач лидарной одометрии.

Вместе с тем классический ICP имеет и системные ограничения. Локальный характер оптимизации делает его чувствительным к начальному приближению: если стартовое положение текущего облака слишком далеко от истинного, поиск ближайших соответствий может начать систематически связывать между собой геометрически несвязанные точки. Кроме того, качество результата напрямую зависит от полноты межкадрового перекрытия и от доли ошибочных пар. Наконец, базовая квадратичная функция ошибки сравнительно слабо подавляет крупные ошибки в отдельных соответствиях, поэтому отдельные неудачные соответствия могут заметно влиять на итоговую оценку преобразования.

Для уменьшения влияния ложных соответствий в задачах регистрации нередко применяются дополнительные процедуры отбора согласованных пар. Одним из наиболее известных инструментов такого типа является \gls{ransac}~\cite{ransac-pyihub}, в котором по случайным подмножествам соответствий многократно строятся гипотезы преобразования, а затем выбирается гипотеза, согласующаяся с наибольшим числом наблюдений. Такой подход особенно полезен в ситуациях, где доля ошибочных пар может быть заметной. Однако даже использование \gls{ransac} не устраняет проблему неполного перекрытия сцены: если информативных статических точек остаётся слишком мало, то после отбрасывания сомнительных соответствий задача оценки движения всё равно может оказаться плохо обусловленной.

Следовательно, ICP удобен как базовый метод сравнения и как отправная точка для построения лидарной одометрии, однако его классическая постановка оказывается недостаточно устойчивой в сценариях, где нарушается предположение о доминировании статической сцены. Именно такая ситуация возникает при динамическом перекрытии поля зрения, которое далее рассматривается отдельно.
